\section{Base points in every Weil family}\label{sec:basepoints}

The secant construction of \Cref{sec:construction} lives on the split variety
$X\times\widehat{X}$ and stops short of the deformation step. In this section
we study the split member and the other special members of a Weil family
through the Weil classes themselves. On the split member the Weil classes are
polynomials in three divisor classes (\Cref{thm:splitclosed}); on products of
abelian surfaces of Weil type they are polynomials in divisor classes
(\Cref{thm:everyn}); and on abelian varieties with quaternionic
multiplication extending $K$, and on their products, they are polynomials in
two divisor classes, which yields an algebraic base point in every Weil
family, of every discriminant (\Cref{thm:everyfamily}). What remains is the
propagation of one flat class over one connected period domain, formulated
in \Cref{ssec:residual,ssec:exactform}.

\subsection{The split member in closed form}\label{ssec:split}

We begin with the split member, where a single formula, valid for every $n$
and every $d$, writes the Weil classes as polynomials in three divisor
classes. The split locus is thin. A general abelian $2n$-fold
of Weil type is simple, hence not isogenous to a product, and for $n\ge2$ the
members $X\times\widehat{X}$ form a subvariety of $\cD_{n,n}$ of positive
codimension (\Cref{prop:splitgeom}(iii)). \Cref{cor:splitalgebraic} is
therefore a base case, and \Cref{q:residual} is a question about propagation.

\begin{notation}\label{not:threeclasses}
On $A=X\times\widehat{X}$ write
\[
  \beta\in H^{2}(A,\QQ),
  \qquad
  \widehat{\beta}\in H^{2}(A,\QQ),
  \qquad
  \ell=c_{1}(\cP)\in H^{2}(A,\QQ)
\]
for the pullback of the polarisation of $X$, the pullback of the dual
polarisation of $\widehat{X}$, and the class of the Poincar\'e bundle. All
three are divisor classes, so every polynomial in them is algebraic. Put
\begin{equation}\label{eq:gammadef}
  \gamma=d\,\beta-\widehat{\beta}.
\end{equation}
\end{notation}

\begin{theorem}[The Weil line of the split member]\label{thm:splitclosed}
Let $A=X\times\widehat{X}$ carry the $K$-action \eqref{eq:Mdef} determined by a
principal $\beta$, and let $V_{\pm}\subset H^{1}(A,\CC)$ be the eigenspaces of
$\iota(\sqrt{-d})$. Then
\begin{enumerate}[label=\textup{(\roman*)},nosep]
\item $\gamma\mp\sqrt{-d}\,\ell$ lies in $\bigwedge^{2}V_{\pm}$, so it is an
  eigenclass for the character $\tau\mapsto\tau^{2}$ on $H^{2}(A,\CC)$
  respectively for its conjugate;
\item the generators of $\bigwedge^{2n}V_{\pm}$ are
  \begin{equation}\label{eq:splitclosed}
    \omega_{\pm}
    =\frac{(-1)^{n(n-1)/2}}{n!\,d^{\,n}}
      \bigl(\gamma\mp\sqrt{-d}\,\ell\bigr)^{n};
  \end{equation}
\item consequently the Weil line of $A$ is
  \[
    \HW(A)=\Bigl\langle\;
      \Re\bigl(\gamma-\sqrt{-d}\,\ell\bigr)^{n},\;\;
      \Im\bigl(\gamma-\sqrt{-d}\,\ell\bigr)^{n}
    \;\Bigr\rangle_{\QQ},
  \]
  and every Weil class on $A$ is a polynomial with rational coefficients in
  the three divisor classes $\beta,\widehat{\beta},\ell$, hence is algebraic.
\end{enumerate}
\end{theorem}

\begin{proof}
Write $U=H^{1}(X,\QQ)$ with basis $x_{1},\dots,x_{2n}$ chosen so that
$\beta=\sum_{j\le n}x_{j}\wedge x_{n+j}$, let $\xi_{1},\dots,\xi_{2n}$ be the
dual basis of $H^{1}(\widehat{X},\QQ)=U^{\vee}$, and let $B\colon U\to U^{\vee}$
be the isomorphism induced by $\beta$, so that $Bx_{j}=\xi_{n+j}$ and
$Bx_{n+j}=-\xi_{j}$. The endomorphism \eqref{eq:Mdef} acts on homology, and
its transpose $M^{*}$ acts on $H^{1}(A,\QQ)=U\oplus U^{\vee}$ by
$u\mapsto -Bu$ and $\phi\mapsto d\,B^{-1}\phi$. Indeed, in the basis
$u_{1},\dots,u_{2n}$ of $H_{1}(X,\QQ)$ dual to the $x_{i}$, the map
$\beta\colon H_{1}(X)\to H^{1}(X)$ sends $u_{j}\mapsto x_{n+j}$ and
$u_{n+j}\mapsto-x_{j}$, so \eqref{eq:Mdef} sends $u_{j}\mapsto dx_{n+j}$,
$u_{n+j}\mapsto-dx_{j}$, $x_{n+j}\mapsto-u_{j}$ and $x_{j}\mapsto u_{n+j}$,
where $x_{i}$ is regarded as an element of
$H_{1}(\widehat X,\QQ)=H^{1}(X,\QQ)$ and $\xi_{i}$ is its dual class; pairing
with the dual basis gives the stated action. The square of $M^{*}$ is $-d$;
so, writing $\delta=\sqrt{-d}$,
\begin{equation}\label{eq:splitVpm}
  V_{\pm}=\Bigl\{\,u\pm\tfrac{\delta}{d}\,Bu\;:\;u\in U_{\CC}\Bigr\},
\end{equation}
since $M^{*}\bigl(u\pm\tfrac{\delta}{d}Bu\bigr)
=-Bu\pm\delta u=\pm\delta\bigl(u\pm\tfrac{\delta}{d}Bu\bigr)$.

(i) Take the sum over $j\le n$ of the wedges of the two eigenvectors attached
to $x_{j}$ and to $x_{n+j}$:
\[
  \sum_{j=1}^{n}
  \Bigl(x_{j}\pm\tfrac{\delta}{d}Bx_{j}\Bigr)\wedge
  \Bigl(x_{n+j}\pm\tfrac{\delta}{d}Bx_{n+j}\Bigr)
  =\beta\pm\tfrac{\delta}{d}\Sigma-\tfrac{1}{d}\widehat{\beta},
\]
using $\delta^{2}=-d$, where
$\Sigma=\sum_{j}\bigl(x_{j}\wedge Bx_{n+j}+Bx_{j}\wedge x_{n+j}\bigr)$ and
$\sum_{j}Bx_{j}\wedge Bx_{n+j}=\sum_{j}\xi_{n+j}\wedge(-\xi_{j})
=\widehat{\beta}$. Now
\[
  \Sigma=\sum_{j=1}^{n}\bigl(x_{j}\wedge(-\xi_{j})
        +\xi_{n+j}\wedge x_{n+j}\bigr)
   =-\sum_{i=1}^{2n}x_{i}\wedge\xi_{i}=-\ell .
\]
So the displayed sum is $\tfrac1d\bigl(d\beta-\widehat{\beta}\mp\delta\ell\bigr)
=\tfrac1d(\gamma\mp\delta\ell)$. It is a sum of wedges of pairs of vectors of
$V_{\pm}$, hence lies in $\bigwedge^{2}V_{\pm}$, which proves (i); the
statement about characters is \Cref{prop:grading} at $(r,s)=(2,0)$ and
$(0,2)$.

(ii) By (i) the $n$-th power $(\gamma\mp\delta\ell)^{n}$ lies in
$\bigwedge^{2n}V_{\pm}$, which is a line by \Cref{prop:weilline}. It is
nonzero: expanding $\gamma\mp\delta\ell=d\beta-\widehat{\beta}\mp\delta\ell$,
the only term of the $n$-th power lying in $\bigwedge^{2n}U$ is
$d^{\,n}\beta^{n}$, and $\beta^{n}\ne0$ because $\beta$ is nondegenerate;
the remaining terms all involve at least one $\xi$, so no cancellation can
occur. Hence $(\gamma\mp\delta\ell)^{n}$ spans $\bigwedge^{2n}V_{\pm}$ and is a
scalar multiple of $\omega_{\pm}=\bigwedge_{i=1}^{2n}\bigl(x_{i}\pm\tfrac{\delta}{d}Bx_{i}\bigr)$.
To find the scalar, compare the components in $\bigwedge^{2n}U$: that of
$\omega_{\pm}$ is $x_{1}\wedge\cdots\wedge x_{2n}$, while
\[
  d^{\,n}\beta^{n}=d^{\,n}\,n!\bigwedge_{j=1}^{n}(x_{j}\wedge x_{n+j})
  =d^{\,n}\,n!\,(-1)^{n(n-1)/2}\,x_{1}\wedge\cdots\wedge x_{2n},
\]
the sign being that of the permutation carrying
$(1,n+1,2,n+2,\dots,n,2n)$ to $(1,2,\dots,2n)$. Inverting gives
\eqref{eq:splitclosed}.

(iii) The two displayed classes are the real and imaginary parts of
$(\gamma-\delta\ell)^{n}$, which up to the nonzero rational factor in
\eqref{eq:splitclosed} are $\omega_{-}+\omega_{+}$ and
$(\omega_{-}-\omega_{+})/\delta$; by \Cref{prop:weilline}(i) these span
$\HW(A)$ over $\QQ$. Expanding $(\gamma-\delta\ell)^{n}$ by the binomial
theorem and separating even from odd powers of $\delta$ writes each of them as
a rational polynomial in $\beta,\widehat{\beta},\ell$, because
$\delta^{2m}=(-d)^{m}$ and $\delta^{2m+1}=(-d)^{m}\delta$. Products of divisor
classes are algebraic, so both are algebraic.
\end{proof}

\begin{example}\label{ex:splitsmall}
The $2n$ classes $x_{1},\dots,x_{2n}$ form a $K$-basis of $H^{1}(A,\QQ)$, on
which $\sqrt{-d}$ acts by $x\mapsto-Bx$, so \eqref{eq:omegaprod} applies to
them; with \eqref{eq:splitclosed} it gives the integral generators of
\Cref{thm:explicitgens} for this basis in closed form,
\[
  \omega_{1}+\sqrt{-d}\,\omega_{2}
  =\frac{(-1)^{n(n-1)/2}}{n!}\,\bigl(\gamma+\sqrt{-d}\,\ell\bigr)^{n}.
\]
At $n=2$ this reads
\[
  \omega_{1}=\tfrac12\bigl(d\,\ell^{2}-\gamma^{2}\bigr),
  \qquad
  \omega_{2}=-\gamma\,\ell,
\]
and at $n=3$,
\[
  \omega_{1}=\tfrac12\,d\,\gamma\,\ell^{2}-\tfrac16\gamma^{3},
  \qquad
  \omega_{2}=\tfrac16\,d\,\ell^{3}-\tfrac12\gamma^{2}\ell,
\]
with $\gamma=d\beta-\widehat{\beta}$ throughout.
\end{example}

\begin{corollary}[Algebraicity on the split locus]\label{cor:splitalgebraic}
For every $n\ge1$, every squarefree $d>0$ and every principally polarised
abelian $n$-fold $X$, the Weil classes of $A=X\times\widehat{X}$ with the
$K$-action \eqref{eq:Mdef} are algebraic, and an algebraic representative is
given by \eqref{eq:splitclosed} as an explicit rational combination of
intersections of three divisors.
\end{corollary}

\begin{proof}
This is \Cref{thm:splitclosed}(iii), together with the Lefschetz theorem on
$(1,1)$ classes, which makes $\beta,\widehat{\beta},\ell$ algebraic, and the
fact that the algebraic classes form a subring of $H^{*}(A,\QQ)$.
\end{proof}

\begin{remark}[Consistency with the character obstruction]\label{rem:splitconsistent}
\Cref{cor:splitalgebraic} does not contradict \Cref{thm:character}. At the
split point $\NS(A)\otimes\QQ$ contains the three independent classes
$\beta,\widehat{\beta},\ell$, so the hypothesis
$\NS(A)\otimes\QQ=\QQ\eta$ of \Cref{prop:notpoly} fails, and this is the
situation of \Cref{rem:special}. What \Cref{thm:character} forbids is a
$\Kx$-eigenclass for the norm character meeting the Weil line, and
$\gamma\mp\sqrt{-d}\,\ell$ is an eigenclass for $\tau^{2}$ and for
$\bbar{\tau}^{2}$, not for the norm. The formula \eqref{eq:splitclosed} is
the first place in the paper where a class in the relevant eigenspace is
exhibited, and such a class exists because the split member has enough
divisors to build one.
\end{remark}

\subsection{Products of abelian surfaces}\label{ssec:products}

A second source of members with enough divisors is the product, an operation
that the Weil line respects. Throughout this subsection $K$ is fixed and all
Weil structures are taken with respect to it.

\begin{lemma}[Products of Weil type]\label{lem:prodweil}
Let $(A_{1},\iota_{1},\eta_{1})$ be of $(K,-1,n_{1})$-Weil type and
$(A_{2},\iota_{2},\eta_{2})$ of $(K,-1,n_{2})$-Weil type. Then
$A=A_{1}\times A_{2}$, with the diagonal action
$\iota=\iota_{1}\times\iota_{2}$ and the product polarisation
$\eta=p_{1}^{*}\eta_{1}+p_{2}^{*}\eta_{2}$, is of
$(K,-1,n_{1}+n_{2})$-Weil type.
\end{lemma}

\begin{proof}
Write $V_{i}=H^{1}(A_{i},\QQ)$, so $V=H^{1}(A,\QQ)=V_{1}\oplus V_{2}$ as
$K$-vector spaces and $\dim_{K}V=2n_{1}+2n_{2}$. The Rosati involution of the
product polarisation is the product of the two Rosati involutions, so it
restricts on $K$ to conjugation, which is the sign $-1$ of
\Cref{def:weiltype}. For the multiplicities, the Hodge decomposition of $V$ is
the direct sum of those of $V_{1}$ and $V_{2}$, and the eigenspace
decomposition is likewise $V_{\pm}=V_{1,\pm}\oplus V_{2,\pm}$. Hence
\[
  \dim_{\CC}V_{+}^{1,0}
  =\dim_{\CC}V_{1,+}^{1,0}+\dim_{\CC}V_{2,+}^{1,0}
  =n_{1}+n_{2},
\]
and the same for the other three, which is \Cref{lem:balanced} for
$n=n_{1}+n_{2}$.
\end{proof}

For an abelian variety $A$ of $(K,-1,n)$-Weil type and a $K$-basis of
$H^{1}(A,\QQ)$, let $\omega_{1}(A),\omega_{2}(A)$ be the generators of the
Weil line given by \eqref{eq:omegagens} in that basis, and write
\[
  \omega(A)=\omega_{1}(A)+\sqrt{-d}\,\omega_{2}(A)\;\in\;
  K\otimes_{\QQ}H^{2n}(A,\QQ)
\]
for the $K$-valued Weil class; by \eqref{eq:omegaprod} its image in
$H^{2n}(A,\CC)$ is $d^{-n}\prod_{j}(d\,x_{j}+\delta\,y_{j})$. On a product we
use the $K$-basis obtained by concatenating bases of the factors.

The Weil line is the top exterior power over $K$, and a top exterior power of
a direct sum is the tensor product of the top exterior powers. This is the
content of the next theorem; what has to be checked is its effect on the
rational generators, which are defined through the real structure rather
than over $K$.

\begin{theorem}[The Weil class is multiplicative]\label{thm:weilmult}
With the notation of \Cref{lem:prodweil}, and cup product on
$H^{*}(A,\QQ)$ taken after pullback along the two projections,
\begin{equation}\label{eq:weilmult}
\begin{aligned}
  \omega_{1}(A)&=\omega_{1}(A_{1})\,\omega_{1}(A_{2})
    -d\,\omega_{2}(A_{1})\,\omega_{2}(A_{2}),\\
  \omega_{2}(A)&=\omega_{1}(A_{1})\,\omega_{2}(A_{2})
    +\omega_{2}(A_{1})\,\omega_{1}(A_{2}),
\end{aligned}
\end{equation}
that is $\omega(A_{1})\,\omega(A_{2})=\omega(A_{1}\times A_{2})$ for the
$K$-valued Weil class. For $k$ factors,
$\prod_{j}\omega(A_{j})=\omega(A_{1}\times\cdots\times A_{k})$.
\end{theorem}

\begin{proof}
Let $u_{1},\dots,u_{2n_{1}}$ and $u_{2n_{1}+1},\dots,u_{2n}$ be the
$K$-bases of $H^{1}(A_{1},\QQ)$ and $H^{1}(A_{2},\QQ)$, so that together they
form the $K$-basis of $H^{1}(A,\QQ)$, and put $x_{j}=u_{j}$ and
$y_{j}=\sqrt{-d}\,u_{j}$ as in \eqref{eq:omegaprod}. By \eqref{eq:omegaprod},
\[
\begin{aligned}
  \omega(A)
  &=d^{-n}\prod_{j=1}^{2n}\bigl(d\,x_{j}+\delta\,y_{j}\bigr)\\
  &=\Bigl(d^{-n_{1}}\prod_{j\le2n_{1}}\bigl(d\,x_{j}+\delta\,y_{j}\bigr)\Bigr)
   \Bigl(d^{-n_{2}}\prod_{j>2n_{1}}\bigl(d\,x_{j}+\delta\,y_{j}\bigr)\Bigr)
  =\omega(A_{1})\,\omega(A_{2})
\end{aligned}
\]
in $H^{2n}(A,\CC)$, the products being taken in the order of the basis, so
that no sign arises. Expanding
$\bigl(\omega_{1}(A_{1})+\delta\omega_{2}(A_{1})\bigr)
\bigl(\omega_{1}(A_{2})+\delta\omega_{2}(A_{2})\bigr)$ with $\delta^{2}=-d$,
and comparing rational parts and $\delta$-parts, which is legitimate because
the classes $\omega_{j}(A_{i})$ are rational and $1,\delta$ are linearly
independent over $\QQ$, gives \eqref{eq:weilmult}. The statement for $k$
factors follows by induction.
\end{proof}

\Cref{fig:product} shows the multiplicativity and the loci it reaches.

\begin{corollary}[Algebraicity propagates along products]\label{cor:prodalg}
If the Weil classes of $A_{1}$ and of $A_{2}$ are algebraic, then so are those
of $A_{1}\times A_{2}$, and an algebraic representative is the rational
combination of external products given by \eqref{eq:weilmult}.
\end{corollary}

\begin{proof}
Algebraic classes are closed under pullback along the projections and under
cup product, and \eqref{eq:weilmult} expresses $\omega_{1}(A)$ and
$\omega_{2}(A)$ in those operations applied to the $\omega_{j}(A_{i})$.
\end{proof}

In dimension two there is nothing to prove.

\begin{lemma}[Dimension two]\label{lem:basecase}
Let $A_{0}$ be of $(K,-1,1)$-Weil type, so $\dim A_{0}=2$. Then
$\HW(A_{0})\subset H^{2}(A_{0},\QQ)$ and both generators are algebraic.
\end{lemma}

\begin{proof}
The Weil classes lie in $H^{2n}=H^{2}$ and are of type $(1,1)$ by
\Cref{prop:weilline}(ii), so they are algebraic by the Lefschetz theorem on
$(1,1)$ classes. No hypothesis on $A_{0}$ beyond the Weil structure is used.
\end{proof}

\begin{theorem}[An algebraic Weil class in every dimension]\label{thm:everyn}
Let $n\ge1$, let $d>0$ be squarefree, and let $A_{1},\dots,A_{n}$ be abelian
surfaces of $(K,-1,1)$-Weil type. Then $A=A_{1}\times\cdots\times A_{n}$ is of
$(K,-1,n)$-Weil type and its Weil classes are algebraic; explicitly, by
\Cref{thm:weilmult} and \Cref{lem:basecase} they are the rational part and
the coefficient of $\sqrt{-d}$ in
$\prod_{j=1}^{n}\bigl(\omega_{1}(A_{j})+\sqrt{-d}\,\omega_{2}(A_{j})\bigr)$,
a polynomial with integer coefficients in $2n$ divisor classes.
\end{theorem}

\begin{proof}
\Cref{lem:prodweil} gives the Weil structure, \Cref{lem:basecase} gives
algebraicity of each factor's classes, and \Cref{cor:prodalg} applied $n-1$
times gives the conclusion; the displayed form is \Cref{thm:weilmult} for $k=n$
factors. Each $\omega_{j}(A_{i})$ is a divisor class by
\Cref{lem:basecase}, so the product is a polynomial in $2n$ of them.
\end{proof}

\begin{remark}[The product locus]\label{rem:prodscope}
\Cref{thm:everyn} is unconditional and holds for every $n$; like
\Cref{cor:splitalgebraic} and \Cref{thm:everyfamily}, it produces an algebraic
Weil class in every dimension using nothing beyond the Lefschetz theorem on
$(1,1)$ classes. The product construction supplies a second family of base
points, defined for every $n$ and every $d$, at which the class is written
down explicitly. It also shows that the algebraicity locus is closed under
products, so that a settled case $n_{0}$ propagates to every multiple of
$n_{0}$, and settled cases propagate to their sums. Combined with Markman's
theorem at $n=2$ \cite{Mar25a,Mar25b}, this gives, for every $n$,
algebraicity on the locus of products of abelian fourfolds and surfaces of
Weil type, which has dimension $2n$ when $n$ is even.

The reach of these results is limited. The family of polarised abelian
$2n$-folds of $(K,-1,n)$-Weil type has dimension $n^{2}$, and the locus of
products of $n$ surfaces has dimension $n$, so for $n\ge2$ it is a proper
closed subvariety, smaller than the quaternionic locus of \Cref{thm:quatweil},
whose dimension is $n(n+1)/2$. A general member of the family is simple,
hence not isogenous to a product, and neither \Cref{thm:everyn} nor
\Cref{cor:splitalgebraic} applies to it. The general member lies outside all
of these loci, and by \Cref{thm:residual} what joins them to it is the
propagation statement.
\end{remark}

\begin{figure}[htbp]
\centering
  \includegraphics[width=0.96\textwidth]{fig_product}
\caption{\Cref{thm:weilmult} and its reach. Left: the Weil line of each factor
is a $K$-line, the Weil line of the product is the tensor product of the two
over $K$, and the $K$-valued Weil class is multiplicative. Algebraicity therefore passes from the factors to the product, and
at $n=1$ the factors are divisor classes, so a product of $n$ abelian surfaces
of Weil type carries an explicit algebraic Weil class in every dimension.
Right: the dimensions of the loci on which the class is known to be algebraic,
against the dimension $n^{2}$ of the family. Both the quaternionic locus and
the product locus are proper subvarieties once $n\ge2$, and a general member
lies on neither.}
\label{fig:product}
\end{figure}

\subsection{The residual statement}\label{ssec:residual}

The results above allow the remaining question to be stated in one sentence,
without auxiliary hypotheses.

\begin{setupenv}[The family]\label{setup:family}
Fix $n\ge2$ and a squarefree $d>0$. Let $\cD_{n,n}$ be the period domain of
\Cref{prop:weildomain} and let
\[
  \pi\colon\cA\longrightarrow\cD_{n,n}
\]
be the family of polarised abelian $2n$-folds of $(K,-1,n)$-Weil type it
parametrises, with $\omega$ the flat section of $R^{2n}\pi_{*}\QQ$ given
fibrewise by the class $\omega_{1}$ of \eqref{eq:omegagens}. Let
$\cD^{\mathrm{spl}}\subset\cD_{n,n}$ be the locus of split members
$X\times\widehat{X}$ carrying the action \eqref{eq:Mdef}.
\end{setupenv}

\begin{question}[The residual statement]\label{q:residual}
Is $\omega_{s}$ algebraic for every $s\in\cD_{n,n}$?
\end{question}

\begin{theorem}[The reduction]\label{thm:residual}
Fix $n\ge2$ and a squarefree $d>0$, and keep the notation of
\Cref{setup:family}.
\begin{enumerate}[label=\textup{(\alph*)},nosep]
\item If \Cref{q:residual} has a positive answer, then the Hodge conjecture
  holds, in every degree, for every abelian $2n$-fold of $(K,-1,n)$-Weil type
  with $\Hg(A)=\SU(V,H)$.
\item Conversely, if the Hodge conjecture holds for every abelian $2n$-fold of
  $(K,-1,n)$-Weil type, then \Cref{q:residual} has a positive answer.
\item $\cD^{\mathrm{spl}}$ is nonempty, $\omega_{s}$ is algebraic at every
  $s\in\cD^{\mathrm{spl}}$ with the explicit representative
  \eqref{eq:splitclosed}, and $\cD_{n,n}$ is connected.
\end{enumerate}
Thus \Cref{q:residual} is the variational Hodge conjecture for one flat class
on one connected family; it is necessary as well as sufficient, and its base
case is supplied. \Cref{thm:everyfamily} supplies a base case in every family,
for every discriminant, so nonemptiness of the base locus is never the
obstacle.
\end{theorem}

\begin{proof}
(a) Let $A$ be of $(K,-1,n)$-Weil type with $\Hg(A)=\SU(V,H)$. By
\Cref{prop:weildomain} its Hodge structure is a point $s$ of $\cD_{n,n}$, so
$A$ is a fibre of $\pi$, and the Weil line of $A$ is generated over $K$ by
$\omega_{s}$. A positive answer makes $\omega_{s}$ algebraic, hence makes the
whole Weil line algebraic by the argument of \Cref{cor:whatisleft}. By
\Cref{thm:hdgcount} the Weil line together with the powers of $\eta$ exhausts
$\Hdg^{*}(A)$, and the powers of $\eta$ are algebraic, so every Hodge class on
$A$ is algebraic.

(b) For every $s$ the class $\omega_{s}$ is a Hodge class on the fibre $A_{s}$,
by \Cref{prop:weilline}(ii), and $A_{s}$ is of $(K,-1,n)$-Weil type by
construction. So the Hodge conjecture for $A_{s}$ gives its algebraicity.

(c) $\cD^{\mathrm{spl}}$ is nonempty because the construction \eqref{eq:Mdef}
applies to any principally polarised abelian $n$-fold and produces a point of
$\cD_{n,n}$ by \Cref{thm:explicitweil}; algebraicity there is
\Cref{cor:splitalgebraic}; and $\cD_{n,n}$ is connected because it is a
bounded symmetric domain, by \Cref{prop:weildomain}.
\end{proof}

\begin{remark}[The target]\label{rem:target}
The preceding sections remove three parts of the problem.
\Cref{thm:hdgcount} settles how many classes are at issue: one, in one
degree, on each fibre. \Cref{thm:divisor} and \Cref{thm:character} exclude
two families of candidate constructions. \Cref{thm:splitclosed} supplies the
base case: the class is algebraic somewhere in the family, with a
representative that can be written on a single line.

The next step is propagation, which is Grothendieck's variational Hodge
conjecture for the pair $(\pi,\omega)$ of \Cref{setup:family}.
\Cref{thm:residual} shows that for the Hodge conjecture on
abelian varieties of Weil type it is necessary as well as sufficient, so every
route to the conjecture meets an equivalent difficulty.
\Cref{ssec:properness} analyses the propagation; there the passage from a
formal deformation to an algebraic one, usually the harder half, is supplied
by properness of the Hilbert scheme. What survives is either an infinitesimal
statement at a single point or a bound on the complexity of representing
cycles along one orbit, and the second is equivalent to its conclusion
(\Cref{thm:p1equivalent}).
\end{remark}

\subsection{The divisorial criterion}
\label{ssec:divcrit}

\Cref{thm:splitclosed} is a computation on one member of the family, but the
mechanism behind it is not special to that member. Isolating it gives a
criterion that can be tested at any point; we then determine where it holds.

\begin{theorem}[The divisorial criterion]\label{thm:divcriterion}
Let $(A,\iota,\eta)$ be of $(K,-1,n)$-Weil type and suppose there exists
$\Theta\in\bigwedge^{2}V_{-}$ with
\begin{enumerate}[label=\textup{(\roman*)},nosep]
\item $\Theta$ nondegenerate as an alternating form on $V_{-}^{\vee}$,
  equivalently $\Theta^{n}\ne0$;
\item $\Theta$ of Hodge type $(1,1)$;
\item $\Theta+\bbar{\Theta}$ rational, that is in $H^{2}(A,\QQ)$.
\end{enumerate}
Then the Weil classes of $A$ are algebraic. Explicitly, writing
$g=\Theta+\bbar{\Theta}$ and $\lambda=(\Theta-\bbar{\Theta})/\sqrt{-d}$, both
$g$ and $\lambda$ are divisor classes and
\begin{equation}\label{eq:divrep}
  \HW(A)=\Bigl\langle\;
    \Re\bigl(g+\sqrt{-d}\,\lambda\bigr)^{n},\;\;
    \Im\bigl(g+\sqrt{-d}\,\lambda\bigr)^{n}\;\Bigr\rangle_{\QQ},
\end{equation}
so every Weil class is a rational polynomial in two divisor classes.
\end{theorem}

\begin{proof}
Conjugation interchanges $\bigwedge^{2}V_{-}$ and $\bigwedge^{2}V_{+}$, so
$\bbar{\Theta}\in\bigwedge^{2}V_{+}$, and both $g$ and $\lambda$ are fixed by
conjugation, hence real. Rationality of $g$ is (iii); that of $\lambda$
comes from the action of $\Kx$. Fix
$\tau=a+b\sqrt{-d}\in\Kx$ with $ab\ne0$. By \Cref{prop:grading} the operator
$\iota(\tau)^{*}$ acts on $\bigwedge^{2}V_{+}$ by $\tau^{2}=a^{2}-db^{2}+2ab\sqrt{-d}$
and on $\bigwedge^{2}V_{-}$ by $\bbar{\tau}^{\,2}$, so
\[
  J:=\frac{\iota(\tau)^{*}-(a^{2}-db^{2})}{2ab}
\]
is a $\QQ$-rational endomorphism of $H^{2}(A,\QQ)$ preserving the space $R$
of \Cref{prop:normpart} and acting
on $\bigwedge^{2}V_{\pm}$ by $\pm\sqrt{-d}$. Applying it to
$g=\Theta+\bbar{\Theta}$ gives $J(g)=-\sqrt{-d}\,\Theta+\sqrt{-d}\,\bbar{\Theta}
=-\sqrt{-d}\cdot\sqrt{-d}\,\lambda=d\,\lambda$, so $\lambda=J(g)/d$ is
rational. Both $g$ and $\lambda$ are of type $(1,1)$ by (ii) and conjugation
symmetry, so both lie in $\NS(A)\otimes\QQ$ by the Lefschetz theorem on
$(1,1)$ classes, and both are algebraic.

By (i) the form $\Theta$ on the $2n$-dimensional space $V_{-}$ has nonzero
Pfaffian, so $\Theta^{n}$ is a nonzero element of the line
$\bigwedge^{2n}V_{-}$, and it therefore spans it; conjugating,
$\bbar{\Theta}^{\,n}$ spans $\bigwedge^{2n}V_{+}$. Hence
$\Theta^{n}+\bbar{\Theta}^{\,n}$ and
$(\Theta^{n}-\bbar{\Theta}^{\,n})/\sqrt{-d}$ span $\HW(A)$ over $\QQ$, by
\Cref{prop:weilline}(i). Finally
$2^{n}\Theta^{n}=(g+\sqrt{-d}\lambda)^{n}$, and expanding by the binomial
theorem with $(\sqrt{-d})^{2m}=(-d)^{m}$ writes the real and imaginary parts
as rational polynomials in $g$ and $\lambda$. Products of divisor classes are
algebraic.
\end{proof}

\begin{corollary}\label{cor:splitsatisfies}
The split member of \Cref{thm:splitclosed} satisfies the criterion, with
$\Theta=\gamma-\sqrt{-d}\,\ell$ and $\gamma=d\beta-\widehat{\beta}$; the two
divisor classes of \eqref{eq:divrep} are $2\gamma$ and $-2\ell$.
\end{corollary}

\begin{proof}
\Cref{thm:splitclosed}(i) places $\Theta=\gamma-\sqrt{-d}\,\ell$ in
$\bigwedge^{2}V_{+}$ and its conjugate in $\bigwedge^{2}V_{-}$. The hypotheses
and the conclusion of \Cref{thm:divcriterion} are symmetric under replacing
$\Theta$ by $\bbar{\Theta}$, which interchanges $V_{+}$ and $V_{-}$, fixes $g$
and changes the sign of $\lambda$, so the criterion may be applied to
$\Theta$. The classes $\gamma$ and $\ell$ are divisor classes, hence of type
$(1,1)$, so $\Theta$ is of type $(1,1)$ and $\Theta+\bbar{\Theta}=2\gamma$ is
rational. Nondegeneracy is the nonvanishing of $\Theta^{n}$, proved in
\Cref{thm:splitclosed}(ii).
\end{proof}

The base case is therefore one instance of a general criterion. The next two
results determine how far the criterion reaches, and both limit it.

\begin{proposition}[The criterion fails at a general point]\label{prop:divfails}
Let $(A,\iota,\eta)$ be of $(K,-1,n)$-Weil type with $n\ge2$ and
$\Hg(A)=\SU(V,H)$. Put
\[
  R=\Bigl(\textstyle\bigwedge^{2}V_{+}\oplus\bigwedge^{2}V_{-}\Bigr)
      \cap H^{2}(A,\QQ).
\]
Then $R\cap\NS(A)_{\QQ}=0$, and no $\Theta$ as in \Cref{thm:divcriterion}
exists.
\end{proposition}

\begin{proof}
By \Cref{prop:normpart} the space $R$ is a rational sub-Hodge structure of
$H^{2}(A,\QQ)$, and the intersection $R\cap\NS(A)_{\QQ}$ is a rational
sub-Hodge structure of $R$ of pure type $(1,1)$. The complexification of a
nonzero rational sub-Hodge structure of $R$ is a nonzero
$\Hg(A)_{\CC}$-subrepresentation of $\bigwedge^{2}V_{+}\oplus\bigwedge^{2}V_{-}$,
so it contains a constituent isomorphic to $\bigwedge^{2}V_{+}$ or to
$\bigwedge^{2}V_{-}$; by the proof of \Cref{prop:normpart}, Hodge numbers
depend only on the representation, and by \Cref{prop:h2pieces} each of
$\bigwedge^{2}V_{\pm}$ has $h^{2,0}+h^{0,2}=2\binom{n}{2}=n(n-1)$, which is
positive for $n\ge2$. So no nonzero rational sub-Hodge structure of $R$ is of
pure type $(1,1)$, and $R\cap\NS(A)_{\QQ}=0$. The argument covers the case
$n=2$ with trivial discriminant, where $R$ is reducible. A $\Theta$ as in
\Cref{thm:divcriterion} would give the nonzero element $\Theta+\bbar{\Theta}$
of $R\cap\NS(A)_{\QQ}$.
\end{proof}

\begin{proposition}[The split locus]\label{prop:splitgeom}
Fix $n\ge1$ and a squarefree $d>0$, and let $A=X\times\widehat{X}$ carry the
$K$-action \eqref{eq:Mdef} determined by a nondegenerate
$\beta\in\NS(X)\otimes\QQ$ of elementary divisors $(d_{1},\dots,d_{n})$.
\begin{enumerate}[label=\textup{(\roman*)},nosep]
\item Let $B$ be the Gram matrix of $\beta$ in a basis of $U$, and let
  $\widehat\beta$ be the form on $U^{\vee}$ with Gram matrix $-B^{-1}$ in the
  dual basis, which for principal $\beta$ is the class of
  \Cref{not:threeclasses}. The polarisation of $A$ is
  $\eta=\beta+\widehat{\beta}/d$, a positive rational multiple of
  $d\beta+\widehat\beta$; its Gram matrix on $H_{1}(A,\QQ)=U\oplus U^{\vee}$
  is the block form $Q=B\oplus(-B^{-1}/d)$, and $M$ is anti-self-adjoint for
  it.
\item In the $K$-basis $u_{1},\dots,u_{2n}$ of $H_{1}(A,\QQ)$ the hermitian
  form of \Cref{lem:hermform} has Gram matrix $\sqrt{-d}\,B$, so
  \[
    \det H=(-d)^{n}\det B=(-d)^{n}(d_{1}\cdots d_{n})^{2},
  \]
  and $(-1)^{n}\det H=d^{\,n}(d_{1}\cdots d_{n})^{2}$ is a norm from $K$.
  Every split member therefore lies in the Weil family of trivial
  discriminant, whatever $\beta$ is chosen.
\item The split members form a locus of dimension $n(n+1)/2$ inside
  $\cD_{n,n}$, of codimension $n(n-1)/2$, which is zero only for $n=1$ and is
  a divisor for $n=2$.
\item For $X$ with $\Hg(X)=\mathrm{Sp}(H^{1}(X,\QQ))$ one has
  $\NS(A)_{\QQ}=\langle\beta,\widehat{\beta},\ell\rangle$, of rank three, and
  in degree $2n$ the monomials in those three classes span the whole space of
  Hodge classes of $A$.
\end{enumerate}
\end{proposition}

\begin{proof}
(i) In the chosen basis of $U$ and the dual basis of $U^{\vee}$ the map
$\beta\colon U\to U^{\vee}$ has matrix $B^{\top}=-B$ and its inverse has
matrix $-B^{-1}$, so \eqref{eq:Mdef} becomes
$M=\begin{psmallmatrix}0&B^{-1}\\-dB&0\end{psmallmatrix}$, with
$M^{2}=-d$ and
$M^{\top}=\begin{psmallmatrix}0&dB\\-B^{-1}&0\end{psmallmatrix}$, using
$B^{\top}=-B$ and $(B^{-1})^{\top}=-B^{-1}$. For a block form
$Q=\begin{psmallmatrix}P&\Lambda\\-\Lambda^{\top}&S\end{psmallmatrix}$ the
Rosati relation $M^{\top}Q+QM=0$ reads
$B\Lambda^{\top}+\Lambda B=0$, $dBS+PB^{-1}=0$, $B^{-1}P+dSB=0$ and
$B^{-1}\Lambda+\Lambda^{\top}B^{-1}=0$. Taking $\Lambda=0$ and $P=B$ forces
$S=-B^{-1}/d$ from the second relation, and the third then reads $-I+I=0$.
So $Q=B\oplus(-B^{-1}/d)$, which is the Gram matrix of
$\beta+\widehat\beta/d$. For principal $\beta$ one may take
$B=\begin{psmallmatrix}0&I\\-I&0\end{psmallmatrix}$, so that $-B^{-1}=B$
and $\widehat\beta=\sum_{j}\xi_{j}\wedge\xi_{n+j}$ as in
\Cref{not:threeclasses}.

(ii) The set $u_{1},\dots,u_{2n}$ is a $K$-basis because
$MU=d\beta(U)=U^{\vee}$, and $H(u_{i},u_{j})=Q(Mu_{i},u_{j})+\sqrt{-d}\,Q(u_{i},u_{j})$
by \eqref{eq:QfromH}; the first term vanishes since $Mu_{i}\in U^{\vee}$,
$u_{j}\in U$ and $Q$ has no cross block by (i), and the second is
$\sqrt{-d}\,\beta_{ij}$. Hence the Gram matrix is $\sqrt{-d}\,B$ and
$\det H=(\sqrt{-d})^{2n}\det B=(-d)^{n}\det B$. The determinant of an
alternating matrix is the square of its Pfaffian, here $(d_{1}\cdots
d_{n})^{2}$. Now $d=\Nm_{K/\QQ}(\sqrt{-d})$ is a norm and every square is a
norm, and the norms form a subgroup of $\QQ^{\times}$, so
$d^{\,n}(d_{1}\cdots d_{n})^{2}$ is a norm.

(iii) A split member is determined by $X$ together with $\beta$; the moduli
space of polarised abelian $n$-folds has dimension $\dim\cA_{n}=n(n+1)/2$,
and the type of $\beta$ is discrete. By \Cref{prop:weildomain} the ambient domain has
dimension $n^{2}$, and $n^{2}-n(n+1)/2=n(n-1)/2$.

(iv) The Hodge classes of $A=X\times\widehat{X}$ are the invariants of
$\Hg(A)=\Hg(X)=\mathrm{Sp}(U)$ acting on
$\bigwedge^{*}(U\oplus U^{\vee})$. In degree two the invariants are the three
lines spanned by $\beta$, $\widehat{\beta}$ and $\ell$, one from each
K\"unneth summand $\bigwedge^{2}U$, $U\otimes U^{\vee}$ and
$\bigwedge^{2}U^{\vee}$. In higher degree, identify $U^{\vee}$ with $U$
through $\beta$, so that $\bigwedge^{*}(U\oplus U^{\vee})\cong
\bigwedge^{*}(U\otimes\QQ^{2})$ as representations of $\mathrm{Sp}(U)$. By the
first fundamental theorem of invariant theory for the symplectic group, the
invariants of $\mathrm{Sp}(U)$ in $U^{\otimes2m}$ are spanned by the tensors
attached to the pairings of the $2m$ factors, each a product of copies of the
invariant element of $U\otimes U$ dual to $\beta$. The projection of
$(U\otimes\QQ^{2})^{\otimes2m}$ onto $\bigwedge^{2m}(U\otimes\QQ^{2})$ is
equivariant and surjective, hence surjective on invariants since
$\mathrm{Sp}(U)$ is reductive, and it carries such a tensor, tensored with
vectors of $\QQ^{2}$, to a product of invariants of degree two. So the
invariants are generated by $\beta$, $\widehat{\beta}$ and $\ell$, and in
degree $2n$ the monomials $\beta^{i}\widehat{\beta}^{j}\ell^{k}$ with
$i+j+k=n$ span the Hodge classes.
\end{proof}

\begin{remark}[Two residual questions]\label{rem:residualmap}
The Weil families for a fixed $K$ and a fixed $n$ are classified by the
discriminant of the hermitian form, an element of the infinite group
$\QQ^{\times}_{>0}/\Nm_{K/\QQ}(K^{\times})$. By \Cref{prop:splitgeom}(ii) the
split construction of \Cref{ssec:split} produces a base point in only one of
them, the family of trivial discriminant; products of abelian surfaces of
Weil type (\Cref{thm:everyn}, \Cref{lem:products}) reach every discriminant,
but only on a locus of dimension $n$. Read together with
\Cref{thm:residual}, this divides \Cref{q:residual} into two questions of
different kinds.
\begin{enumerate}[label=\textup{(R\arabic*)},leftmargin=2.8em,itemsep=2pt]
\item \emph{Trivial discriminant.} A base point exists, with the explicit
  cycle \eqref{eq:splitclosed}, and what is missing is propagation off a locus
  of codimension $n(n-1)/2$. At $n=2$ that locus is a divisor and the
  propagation is known, by Markman \cite{Mar25a,Mar25b} and by Floccari and Fu
  \cite{FF26}; combined with \Cref{thm:residual}(a) this recovers the Hodge
  conjecture for abelian fourfolds of Weil type of trivial discriminant with
  full Hodge group. At $n=3$ it is settled as well, by Markman's theorem on
  the sixfolds of trivial discriminant \cite{Mar25a,Mar25b}, for every $K$.
  For $n\ge4$ it is the propagation statement, and by \Cref{prop:descent}
  it is then the only question in the imaginary quadratic case.
\item \emph{Nontrivial discriminant.} The split construction never lands
  there, by \Cref{prop:splitgeom}(ii), so it supplies no base
  point; and by \Cref{prop:divfails} no divisorial argument can supply one at
  a general member. A base point must come from a special member at which $R$
  acquires Hodge classes.
\end{enumerate}
Question (R2) is settled in \Cref{ssec:everyfamily}: the special members
required are built from abelian varieties with quaternionic multiplication
extending $K$ and their products, they exist in every family, and on them
\Cref{thm:divcriterion} applies. What then remains is the propagation
question of (R1), posed in every family, each of which has a base point.
\end{remark}

\begin{remark}[Codimension of the split locus]
\label{rem:nlcodim}
A component of the Noether--Lefschetz locus of a weight two variation has
codimension at most $h^{2,0}$, here $n(n-1)$. The split locus has codimension
$n(n-1)/2$ by \Cref{prop:splitgeom}(iii), half of that, so it is a
component of half the expected codimension. This reflects
\Cref{prop:splitgeom}(iv): at a split point $R$ acquires two independent Hodge
classes at once, namely $\gamma$ and $\ell$, and each of them costs
$h^{2,0}/2$ conditions rather than $h^{2,0}$. Whether there are components
of the expected codimension carrying only one class, and whether the
criterion of \Cref{thm:divcriterion} can be met on them, is not needed
below: question (R2) is answered in \Cref{ssec:everyfamily}
through quaternionic multiplication.
\end{remark}

\subsection{A base point in every family}\label{ssec:everyfamily}

By \Cref{prop:splitgeom}(ii) the split construction reaches a single
discriminant class. \Cref{thm:divcriterion} does not depend on the origin of
$\Theta$, and in this subsection $\Theta$ comes from an endomorphism rather
than from a pair of divisors chosen by hand. This source reaches every class,
and it settles question (R2) of \Cref{rem:residualmap}.

\begin{definition}[Discriminant]\label{def:disc}
For $(A,\iota,\eta)$ of $(K,-1,n)$-Weil type with hermitian form $H$ as in
\Cref{lem:hermform}, the \emph{discriminant} is the class of
$(-1)^{n}\det H$ in $\QQ^{\times}_{>0}/\Nm_{K/\QQ}(K^{\times})$. The sign
makes the representative positive, because $H$ has signature $(n,n)$; and the
class is well defined because changing the $K$-basis multiplies $\det H$ by a
norm. Hermitian forms of rank $2n$ over $K$ with signature $(n,n)$ are
classified by this class, so it labels the Weil families.
\end{definition}

The classes that enter \Cref{thm:divcriterion} have a description by
endomorphisms: a class in $R\cap\NS(A)_{\QQ}$ is the same thing as a
symmetric endomorphism that conjugates $K$.

\begin{lemma}[Classes in $R$ and endomorphisms]\label{lem:Rendomorphism}
Let $(A,\iota,\eta)$ be of $(K,-1,n)$-Weil type. The map $\mu\mapsto
g_{\mu}$, $g_{\mu}(x,y)=\eta(\mu x,y)$, is a bijection between
\begin{enumerate}[label=\textup{(\alph*)},nosep]
\item the elements $\mu\in\End^{0}(A)$ with $\mu^{\dagger}=\mu$ and
  $\mu\,\iota(\tau)=\iota(\bbar{\tau})\,\mu$ for all $\tau\in K$, and
\item the elements of $R\cap\NS(A)_{\QQ}$.
\end{enumerate}
Under it, $\mu$ invertible corresponds to $g_{\mu}$ nondegenerate, hence to a
nondegenerate $\Theta$. If such a $\mu\ne0$ exists and the centraliser of
$\iota(K)$ in $\End^{0}(A)$ is $\iota(K)$, then $\mu^{2}\in\QQ^{\times}$ and
$\iota(K)+\iota(K)\mu$ is a quaternion algebra $(-d,\mu^{2})_{\QQ}$ inside
$\End^{0}(A)$.
\end{lemma}

\begin{proof}
A class in $\NS(A)_{\QQ}$ is a symmetric homomorphism $A\to\Av$ up to isogeny,
that is $\eta\mu$ with $\mu^{\dagger}=\mu$. The condition
$g_{\mu}(\iota(\tau)x,y)=g_{\mu}(x,\iota(\tau)y)$, which says that $g_{\mu}$ is
$K$-bilinear and hence lies in $R=\bigwedge^{2}_{K}V$, reads
$\eta(\mu\iota(\tau)x,y)=\eta(\mu x,\iota(\tau)y)=\eta(\iota(\tau)^{\dagger}\mu
x,y)=\eta(\iota(\bbar{\tau})\mu x,y)$, using \Cref{def:weiltype}(iii), so it
is the condition $\mu\iota(\tau)=\iota(\bbar{\tau})\mu$. For the last sentence,
$\mu^{2}$ commutes with $\iota(K)$, hence lies in $\iota(K)$ by hypothesis, and
$\mu\mu^{2}\mu^{-1}=\mu^{2}$ forces it to be fixed by conjugation, so
$\mu^{2}\in\QQ^{\times}$. Then $\iota(\sqrt{-d})$ and $\mu$ generate the
quaternion algebra with the stated invariants.
\end{proof}

The criterion therefore asks for quaternionic multiplication extending $K$.
Abelian varieties with this structure are easy to construct, and their
invariants can be computed.

\begin{setupenv}[The quaternionic model]\label{setup:quat}
Let $b\in\QQ_{>0}$ and let $B=(-d,b)_{\QQ}$ be the quaternion algebra with
basis $1,i,j,ij$ and $i^{2}=-d$, $j^{2}=b$, $ij=-ji$. Let
$V=B^{n}$ as a left $B$-module, of dimension $4n$ over $\QQ$, let
$K=\QQ(i)\subset B$ act through $i$, and fix $a_{1},\dots,a_{n}\in\QQ^{\times}$
and
\begin{equation}\label{eq:quatE}
  E(x,y)=\sum_{k=1}^{n}a_{k}\,\mathrm{trd}\bigl(\bbar{x_{k}}\,i\,y_{k}\bigr),
\end{equation}
where $\mathrm{trd}$ is the reduced trace and $\bbar{\phantom{x}}$ the
canonical involution of $B$.
\end{setupenv}

\begin{theorem}[Weil type from quaternionic multiplication]\label{thm:quatweil}
With the data of \Cref{setup:quat}:
\begin{enumerate}[label=\textup{(\roman*)},nosep]
\item $E$ is alternating and nondegenerate, and its Rosati involution is
  $\beta\mapsto i^{-1}\bbar{\beta}i$, which restricts to complex conjugation on
  $K$ and fixes $j$; it is positive, the form $\mathrm{trd}(\beta\beta^{\dagger})$
  being diagonal with entries $2,2d,2b,2db$ on the basis;
\item in the $K$-basis $e_{1},e_{1}j,\dots,e_{n},e_{n}j$ of $V$ the hermitian
  form of \Cref{lem:hermform} is
  \[
    H=\mathrm{diag}\bigl(2a_{1}d,\,-2a_{1}bd,\;\dots,\;2a_{n}d,\,-2a_{n}bd\bigr),
  \]
  of signature $(n,n)$;
\item there is a complex structure on $V_{\RR}$ commuting with $B$ and
  polarised by $E$, and for every such the resulting abelian variety $A$ is of
  $(K,-1,n)$-Weil type, the balanced signature being forced rather than
  imposed;
\item the discriminant of $A$ is the class of $b^{n}$;
\item $\mu=$ left multiplication by $j$ satisfies \Cref{lem:Rendomorphism}(a),
  so $g=g_{\mu}$ lies in $R\cap\NS(A)_{\QQ}$ and is nondegenerate, the
  divisorial criterion \Cref{thm:divcriterion} applies, and the Weil classes
  of $A$ are algebraic.
\end{enumerate}
\end{theorem}

\begin{proof}
(i) $\mathrm{trd}(z)=\mathrm{trd}(\bbar{z})$ and
$\overline{\bbar{x}iy}=\bbar{y}\,\bbar{i}\,x=-\bbar{y}ix$ give
$E(y,x)=-E(x,y)$. Nondegeneracy is the nondegeneracy of the reduced trace
pairing. For the Rosati, $E(\beta x,y)=E(x,\beta^{\dagger}y)$ for all $x,y$
reads $\bbar{\beta}i=i\beta^{\dagger}$, that is
$\beta^{\dagger}=i^{-1}\bbar{\beta}i$. Then $i^{\dagger}=-i$ and, since
$(-j)i=ij$, also $j^{\dagger}=i^{-1}(ij)=j$. Positivity is the displayed
diagonal, which is positive because $d>0$ and $b>0$.

(ii) $V=Ke_{k}\oplus Ke_{k}j$ over all $k$, since $B=K\oplus Kj$. By
\eqref{eq:QfromH}, $H(x,y)=E(ix,y)+\sqrt{-d}\,E(x,y)$. On one factor,
$E(ix,y)=a\,\mathrm{trd}(\overline{ix}\,i\,y)=-a\,\mathrm{trd}(\bbar{x}i^{2}y)
=ad\,\mathrm{trd}(\bbar{x}y)$, so
\[
  H(x,y)=a\bigl[d\,\mathrm{trd}(\bbar{x}y)+\sqrt{-d}\,\mathrm{trd}(\bbar{x}iy)\bigr].
\]
Now $\mathrm{trd}(1)=2$, $\mathrm{trd}(i)=\mathrm{trd}(j)=\mathrm{trd}(ij)=0$,
$\bbar{j}j=-b$ and $\bbar{j}ij=ib$, which give $H(1,1)=2ad$,
$H(j,j)=-2abd$ and $H(1,j)=0$. Since $b>0$ the two entries of each block have
opposite signs, so the signature is $(n,n)$.

(iii) A complex structure commuting with $B$ and compatible with $E$ is the
same as a subspace $V^{1,0}\subset V_{\CC}$, stable under $B$, with
$V_{\CC}=V^{1,0}\oplus\overline{V^{1,0}}$ and the two Riemann relations. Being
$K$-stable it splits as $P\oplus P'$ with $P\subset V_{+}$ and
$P'\subset V_{-}$; left multiplication by $j$ is $\sigma$-semilinear, so it
carries $V_{+}$ to $V_{-}$, and $B$-stability forces $P'=jP$. Conversely, for
any $n$-dimensional $H$-negative definite $P\subset V_{+}$ the subspace
$P\oplus jP$ is $B$-stable and satisfies the Riemann relations, by the
computation in the proof of \Cref{thm:model} applied blockwise; such $P$
exists because $H$ has signature $(n,n)$ by (ii). Since $j$ is invertible and
carries $P$ isomorphically to $P'$, one gets $\dim P=\dim P'=n$, which is
\Cref{def:weiltype}(ii). The general theory of abelian varieties with
prescribed endomorphism structure is in \cite[Ch.~9]{BL04}.

(iv) By (ii), $\det H=\prod_{k}(2a_{k}d)(-2a_{k}bd)=(-1)^{n}b^{n}\bigl(2^{n}d^{n}
\prod_{k}a_{k}\bigr)^{2}$, so $(-1)^{n}\det H$ is $b^{n}$ times a square, and
squares are norms.

(v) Left multiplication by $j$ satisfies $\mu^{\dagger}=\mu$ by (i) and
$\mu\,\iota(\tau)=\iota(\bbar{\tau})\,\mu$ because $ji=-ij$; it is invertible
because $\mu^{2}=b\ne0$. \Cref{lem:Rendomorphism} puts $g_{\mu}$ in
$R\cap\NS(A)_{\QQ}$ and makes it nondegenerate, and
\Cref{thm:divcriterion} applies to the component $\Theta$ of $g_{\mu}$ in
$\bigwedge^{2}V_{-}$.
\end{proof}

\begin{lemma}[Products and discriminants]\label{lem:products}
Let $(A_{1},\iota_{1},\eta_{1})$ and $(A_{2},\iota_{2},\eta_{2})$ be of
$(K,-1,n_{1})$- and $(K,-1,n_{2})$-Weil type. Then
$A_{1}\times A_{2}$ with the diagonal $K$-action and the product polarisation
is of $(K,-1,n_{1}+n_{2})$-Weil type, its hermitian form is the orthogonal
sum, its discriminant is the product of the two discriminants, and if both
factors satisfy the divisorial criterion then so does the product.
\end{lemma}

\begin{proof}
The first assertion is \Cref{lem:prodweil}. Signatures and hermitian forms add
over an orthogonal sum, so the signature is
$(n_{1}+n_{2},n_{1}+n_{2})$ and $\det H=\det H_{1}\det H_{2}$, whence
$(-1)^{n_{1}+n_{2}}\det H$ is the product of the two normalised determinants.
For the criterion, $V_{-}=V_{1,-}\oplus V_{2,-}$, and
$\Theta=\Theta_{1}+\Theta_{2}$ lies in $\bigwedge^{2}V_{-}$, is of type
$(1,1)$, has $\Theta+\bbar{\Theta}$ rational, and is nondegenerate on the
direct sum because each summand is nondegenerate on its factor.
\end{proof}

\begin{theorem}[A base point in every Weil family]\label{thm:everyfamily}
Let $K=\QQ(\sqrt{-d})$, let $n\ge1$, and let $\delta$ be any class in
$\QQ^{\times}_{>0}/\Nm_{K/\QQ}(K^{\times})$. Then there is an abelian
$2n$-fold $A$ of $(K,-1,n)$-Weil type with discriminant $\delta$ whose Weil
classes are algebraic, with an explicit representative as a rational
polynomial in two divisor classes.
\end{theorem}

\begin{proof}
Choose $b\in\QQ_{>0}$ representing $\delta$.

If $n$ is odd, take the quaternionic model of \Cref{setup:quat} with this $b$.
Its discriminant is $b^{n}=b\cdot(b^{(n-1)/2})^{2}$, which is $\delta$, and
its Weil classes are algebraic by \Cref{thm:quatweil}(v).

If $n$ is even, then $b^{n}$ is a square and the model alone reaches only the
trivial class. Take instead $A=A_{1}\times A_{2}$ with $A_{1}$ the
quaternionic model at $n_{1}=1$ and parameter $b$, and $A_{2}$ the
quaternionic model at $n_{2}=n-1$ and parameter $d$. Both factors satisfy the
criterion, so the product does by \Cref{lem:products}; and the discriminant is
$b^{1}\cdot d^{\,n-1}$, which is $\delta$ because $d=\Nm_{K/\QQ}(\sqrt{-d})$ is
a norm. The dimension is $2(1+(n-1))=2n$.

In both cases the representative is \eqref{eq:divrep} for the class
$g_{\mu}$ of \Cref{thm:quatweil}(v), or, in the second case, for the sum of
the two such classes.
\end{proof}

\begin{corollary}[Reduction to propagation]\label{cor:R2}
Every Weil family contains a member at which the Weil classes are algebraic.
Consequently \Cref{q:residual} reduces to propagation alone: the Hodge
conjecture for abelian $2n$-folds of Weil type, for every $n$ and every
discriminant, holds if and only if the class $\omega$ of
\Cref{setup:family} propagates from the nonempty locus where it is known to be
algebraic to the rest of its connected family.
\end{corollary}

\begin{proof}
The first sentence is \Cref{thm:everyfamily} applied to each class, together
with the fact that the families are indexed by those classes
(\Cref{def:disc}). The second is \Cref{thm:residual}, the base locus being
nonempty in every family by the first.
\end{proof}

\begin{remark}[Two features of the residual statement]\label{rem:whatsleftnow}
By \Cref{thm:residual} and \Cref{cor:R2}, the residual statement is the
variational Hodge conjecture for one flat class on one connected family, with
a base point in every family, and it is necessary as well as sufficient. Two
of its features show which tools cannot settle it.

The base locus is thin. By \Cref{prop:divfails} the divisorial criterion fails
at every point with full Hodge group, so the locus where it holds is contained
in the Noether--Lefschetz locus of $R$, a countable union of proper closed
algebraic subvarieties of $\cD_{n,n}$ by \Cref{thm:cdk}. The divisorial
method cannot enlarge the base locus beyond special points: every point it
produces is again special.

The propagation is not formal. By \Cref{rem:variational} its infinitesimal
form would follow from semiregularity, in the sense of Bloch, of the cycle at
a base point, and that cycle is explicit in every family, being a rational
polynomial in two divisor classes. Its semiregularity is a computation in the
normal bundle that we do not carry out; for the perfect complexes of
\Cref{prop:ktheory}, whose Chern character is concentrated in the degree of
the Weil class, semiregularity fails (\Cref{prop:concentrated},
\Cref{thm:p2false}). The other half of the difficulty, that semiregularity
gives deformations only in the formal category, is removed in
\Cref{ssec:properness}: once the algebraic locus is known to be a countable
union of closed subvarieties, a formal deformation at one point is enough.
\end{remark}

The secant route of \Cref{sec:construction} needs, at $n\ge4$, an object on
$X$ whose Chern character lies on the secant plane and which also supplies the
deformation step of \Cref{rem:deformation}. No such object is known, but its
numerical invariants raise no obstruction.

\begin{proposition}[Numerical admissibility of the secant plane]%
\label{prop:notforbidden}
Let $t$ be a principal polarisation on $X$ and let
$\ch=au+bv$ be a rational point of the secant plane of \Cref{thm:plane},
with $a,b\in\ZZ$. Then
\begin{enumerate}[label=\textup{(\roman*)},nosep]
\item every $\ch_{k}$ is an integral multiple of the integral class
  $t^{k}/k!$, the multiple being $\pm a d^{j}$ in even degree $k=2j$ and
  $\pm b d^{j}$ in odd degree $k=2j+1$;
\item for all integers $a$, $b$ and every positive integer $d$, the Chern
  classes obtained from $\ch$ by Newton's identities are integral in every
  degree: $c_{k}=\gamma_{k}\,t^{k}/k!$ with $\gamma_{0}=1$, $\gamma_{1}=b$ and
  \[
    \gamma_{k+1}\;=\;b\,\gamma_{k}+d\,k\,(a-k+1)\,\gamma_{k-1}\qquad(k\ge1);
  \]
\item the discriminant is
  $\Delta=2rc_{2}-(r-1)c_{1}^{2}=\bigl(a^{2}d+b^{2}\bigr)t^{2}
  =\Nm_{K/\QQ}(b+a\sqrt{-d})\,t^{2}$, which is positive, so Bogomolov's
  inequality never binds.
\end{enumerate}
\end{proposition}

\begin{proof}
(i) is the definition of $u$ and $v$: the coefficient of $t^{k}$ is
$a(-d)^{j}/(2j)!$ or $b(-d)^{j}/(2j+1)!$, and multiplying by $k!$ clears the
factorial and leaves $\pm ad^{j}$ or $\pm bd^{j}$. For (ii), write the total
Chern class as a power series $c(t)=\sum_{k}\gamma_{k}t^{k}/k!$ in the
variable $t$. Newton's identities in generating-function form say
$\log c=\sum_{k\ge1}(-1)^{k-1}(k-1)!\,\ch_{k}$, and with the coefficients of
(i) the two halves of this sum are recognised series:
\[
  \log c(t)\;=\;-a\sum_{j\ge1}\frac{(-d)^{j}t^{2j}}{2j}
     +b\sum_{j\ge0}\frac{(-d)^{j}t^{2j+1}}{2j+1}
  \;=\;\frac{a}{2}\log\bigl(1+dt^{2}\bigr)
     +\frac{b}{\sqrt{d}}\arctan\bigl(\sqrt{d}\,t\bigr).
\]
Differentiating, $(1+dt^{2})\,c'(t)=(b+adt)\,c(t)$. Comparing the
coefficients of $t^{k}/k!$ on the two sides gives
$\gamma_{k+1}+d\,k(k-1)\gamma_{k-1}=b\gamma_{k}+ad\,k\,\gamma_{k-1}$, which is
the stated recursion. Its coefficients are integers and it starts from
$\gamma_{0}=1$, $\gamma_{1}=b$, so every $\gamma_{k}$ is an integer, and
$c_{k}=\gamma_{k}\,t^{k}/k!$ is an integral class because $t^{k}/k!$ is. (iii) is the computation recorded in
\Cref{thm:secantinvariants}, and the norm form of an imaginary quadratic field
is positive definite.
\end{proof}

\begin{remark}[Numerical consistency of the secant object]%
\label{rem:notforbidden}
\Cref{prop:notforbidden} excludes a purely arithmetic obstruction. If the
object required by the construction at $n\ge4$ were forbidden by its
numerical invariants, the secant route would fail on arithmetic grounds
alone. It is not: the Chern character is integral, the Chern classes are
integral in every degree by \Cref{prop:notforbidden}(ii), the
discriminant is positive, and neither Riemann--Roch nor Bogomolov's
inequality stands in the way.
Together with \Cref{prop:objectsize}, which bounds the size of the object from
below without excluding it, and the parameter count of \Cref{rem:markmancheck},
which fixes the rank one shape at $n=3$, every count available to us is
consistent with the existence of the object, and none of them produces it.
What is missing at $n\ge4$ is the object itself.
\end{remark}

\subsection{The quaternionic locus as a Lagrangian Grassmannian}
\label{ssec:lagrangian}

\Cref{thm:quatweil} produces points of the family carrying an admissible
endomorphism $\psi$. The locus of all such points has dimension $n(n+1)/2$,
by a Morita argument recalled in \Cref{rem:ntwo}. Here we identify the locus
itself, and the identification shows where the residual difficulty lies.

\begin{proposition}[The quaternionic locus]\label{prop:lagrangian}
Keep \Cref{setup:quat} and let $\psi$ be $K$-semilinear, invertible,
Rosati-symmetric, with $\psi^{2}=b$. For $x,y\in V_{+}$ put
\[
  \omega_{\psi}(x,y)\;=\;E(x,\psi y).
\]
\begin{enumerate}[label=\textup{(\roman*)},nosep]
\item $\omega_{\psi}$ is a nondegenerate alternating form on $V_{+}$, and it
  is alternating precisely because $\psi$ is Rosati-symmetric.
\item A point $P=V_{+}^{1,0}$ of $\cD_{n,n}$ lies on the locus where $\psi$
  is a morphism of Hodge structures if and only if $P$ is isotropic for
  $\omega_{\psi}$.
\item Since $\dim_{\CC}V_{+}=2n$ and $\dim_{\CC}P=n$, such a $P$ is
  Lagrangian, so the locus is the open subset of
  $\mathrm{LG}(n,V_{+},\omega_{\psi})$ cut out by negative definiteness of
  $H$. Its dimension is $n(n+1)/2$ and its codimension in $\cD_{n,n}$ is
  $n(n-1)/2$.
\end{enumerate}
\end{proposition}

\begin{proof}
(i) $\psi$ carries $V_{+}$ to $V_{-}$, since it is semilinear, and $E$ pairs
$V_{+}$ with $V_{-}$ perfectly, so $\omega_{\psi}$ is defined and
nondegenerate. Rosati symmetry is the identity $E(\psi x,y)=E(x,\psi y)$;
with $E$ alternating it gives
$\omega_{\psi}(y,x)=E(y,\psi x)=-E(\psi x,y)=-E(x,\psi y)=-\omega_{\psi}(x,y)$.

(ii) Suppose $\psi$ preserves the Hodge structure. Being semilinear it
carries $V_{+}^{1,0}$ into $V_{-}^{1,0}$, so $\psi P\subseteq V^{1,0}$, and
$E$ vanishes on $V^{1,0}$ by the first Riemann relation; hence
$E(P,\psi P)=0$, which is the isotropy of $P$. Conversely, if $P$ is
isotropic then $\psi P$ lies in the $E$-annihilator of $P$ inside $V_{-}$.
That annihilator has dimension $n$, because $E$ pairs $V_{+}$ with $V_{-}$
perfectly, and it contains $V_{-}^{1,0}$, which also has dimension $n$ and is
$E$-orthogonal to $P$ since both lie in $V^{1,0}$. The two therefore agree,
so $\psi P=V_{-}^{1,0}$ and $\psi$ preserves $V^{1,0}$.

(iii) A Lagrangian Grassmannian $\mathrm{LG}(n,2n)$ has dimension $n(n+1)/2$,
its tangent space at a Lagrangian $P$ being the symmetric forms on $P$;
negative definiteness is an open condition and does not change the dimension.
Subtracting from $\dim\cD_{n,n}=n^{2}$ gives $n(n-1)/2$.
\end{proof}

\Cref{fig:lagsweep} draws these loci as a pencil of leaves.

\begin{figure}[htbp]
\centering
  \includegraphics[width=0.80\textwidth]{fig_lagsweep}
\caption{The quaternionic loci of \Cref{prop:lagrangian}, drawn as a
pencil. Each admissible $\psi$ gives one leaf, an open piece of
$\mathrm{LG}(n,V_{+},\omega_{\psi})$ of dimension $n(n+1)/2$. As $\psi$
varies over the reals the leaves through a common axis form a continuous
family and sweep out $\cD_{n,n}$. The two heavy leaves are those with $\psi$
rational, and only such leaves are quaternionic loci of abelian varieties.
The drawing is faithful to the continuity of the family and to the
discreteness of its rational members, but not to dimensions: in the family
the leaves have dimension $n(n+1)/2$ inside $n^{2}$, hence codimension
$n(n-1)/2$, which vanishes only at $n=1$ (at $n=2$ the leaves are divisors,
\Cref{rem:ntwo}), and $\psi$ ranges over a real variety rather than an
interval.}
\label{fig:lagsweep}
\end{figure}

\begin{remark}[Real and rational quaternionic loci]%
\label{rem:onlyrational}
Fix a point $P$ of $\cD_{n,n}$, and let $\psi_{0}$ be left multiplication by
$j$ at a point $P_{0}$ given by \Cref{thm:quatweil}(iii). The group
$\SU(V,H)(\RR)$ acts transitively on the bounded symmetric domain $\cD_{n,n}$
of \Cref{prop:weildomain}, so $P=gP_{0}$ for some real $g$; since $g$ is
$K$-linear and preserves $H$, hence $E$, the map
$\psi=g\psi_{0}g^{-1}$ is $K$-semilinear, invertible and Rosati-symmetric,
satisfies $\psi^{2}=b$, and preserves the Hodge structure at $P$. It
satisfies every condition of \Cref{prop:lagrangian} except rationality, and
$P$ lies on its locus. So the quaternionic loci, taken over the \emph{real}
$\psi$, cover $\cD_{n,n}$; they cover it with room to spare, since each has
dimension $n(n+1)/2$ and they are parametrised by a positive-dimensional
family. What prevents the quaternionic loci of abelian varieties from
covering it is that $\psi$ must lie in $\End^{0}(A)$, hence be rational, and
the rational $\psi$ form a countable set.

This is the dichotomy of \Cref{cor:allornothing} in its simplest form. The
obstruction is the countability of a set of parameters inside an uncountable
one, and it involves neither deformation theory nor positivity. What is
needed is a mechanism that produces, from one rational $\psi$, a family of
cycles representing $\omega$ over an open subset of $\cD_{n,n}$ rather than
over the locus of $\psi$ alone; by \Cref{cor:allornothing}, such a family
would answer \Cref{q:residual} in the affirmative.
\end{remark}

\subsection{Absolutely Hodge classes and the exact form}\label{ssec:exactform}

The algebraicity of a Hodge class is usually approached through an
intermediate property. A rational class of type $(p,p)$ on a smooth
projective variety $X$ over $\CC$ is \emph{absolutely Hodge} if, for every
automorphism $\sigma$ of $\CC$, the conjugate under $\sigma$ of its de Rham
component is the de Rham component of a rational class of type $(p,p)$ on the
conjugate variety $X^{\sigma}$. The implications
\begin{equation}\label{eq:threechain}
  \text{algebraic}
  \;\Longrightarrow\;
  \text{absolutely Hodge}
  \;\Longrightarrow\;
  \text{Hodge}
\end{equation}
are both elementary. The Hodge conjecture asserts the converse of their
composite. The converse of the second is the weaker statement, and Deligne
proved it for abelian varieties, so on abelian varieties the Hodge conjecture
is the converse of the first.

\begin{theorem}[Deligne \cite{Del82}]\label{thm:deligneAH}
Every Hodge class on a complex abelian variety is absolutely Hodge. Moreover,
\emph{Principle B}: if $\pi\colon\cX\to S$ is smooth projective over a
connected smooth base, and $\alpha$ is a flat section of $R^{2p}\pi_{*}\QQ(p)$
that is Hodge at every point of $S$ and absolutely Hodge at one point, then it
is absolutely Hodge at every point.
\end{theorem}

Together with \Cref{thm:residual} and \Cref{thm:everyfamily}, this locates
the remaining problem.

\begin{corollary}[Reduction to an absolutely Hodge class]\label{cor:exactform}
Let $\pi\colon\cA\to\cD_{n,n}$ and $\omega$ be as in \Cref{setup:family}. Then
$\omega_{s}$ is absolutely Hodge for every $s$, and the following are
equivalent:
\begin{enumerate}[label=\textup{(\roman*)},nosep]
\item the Hodge conjecture holds, in every degree, for every abelian
  $2n$-fold of $(K,-1,n)$-Weil type with $\Hg(A)=\SU(V,H)$;
\item the absolutely Hodge class $\omega_{s}$ is algebraic for every $s$.
\end{enumerate}
This holds for every $K$, every $n$ and every discriminant.
\end{corollary}

\begin{proof}
The first assertion is \Cref{thm:deligneAH}. It also follows from Principle B
together with \Cref{thm:everyfamily}, which supplies a point of $\cD_{n,n}$ at
which $\omega$ is algebraic and hence absolutely Hodge, and with the
connectedness of $\cD_{n,n}$ from \Cref{prop:weildomain}.

That (ii) implies (i) is \Cref{thm:residual}(a). Conversely, the Hodge group
of $A_{s}$ is $\SU(V,H)$ for $s$ outside a countable union of proper closed
analytic subsets of $\cD_{n,n}$, by Weil's theorem \cite{Weil77} recalled
before \Cref{thm:hdgcount}. Under (i) the class $\omega_{s}$ is algebraic at
all such $s$, so the set $\Sigma$ of points at which it is algebraic is not
meagre, and \Cref{cor:allornothing} gives $\Sigma=\cD_{n,n}$, which is (ii).
\end{proof}
